\chapter[Supercombinators and Lambda-Lifting]{Supercombinators and\\Lambda-Lifting}
%\vspace{2cm}
\vspace{1cm}

\noindent
Since the operation of constructing an instance of a lambda body while
substituting for the formal parameter is the fundamental operation of our
implementation, we will now consider how to make it more efficient.

In this chapter and the next we will show how to transform a lambda
expression into a form in which the lambda abstractions are particularly easy
to instantiate. These special lambda abstractions are called
\textit{supercombinators}, and the transformation is called \textit{lambda-lifting}. Then, in
Chapter~15, we will show how to enhance the lambda-lifting transformation to
be \textit{fully lazy}, a property alluded to in Section~12.1.3. The terms
`supercombinator' and `fully lazy' were both coined by Hughes, who was the
first to combine full laziness with lambda-lifting [Hughes, 1984].

\section{The Idea of Compilation}

The operation of instantiating the body of a lambda abstraction was called
\ml{Instantiate} in the previous chapter, and was performed by a recursive tree-walk
over the lambda body. Such an \ml{Instantiate} operation is rather inefficient for the
following reasons:

\begin{numbered}
    \item at each node of the body, \ml{Instantiate} has to do a case analysis on the tag of
    the node;
    \item at each variable node \ml{Instantiate} has to test if the node is the formal
    parameter; a similar test has to be made at each lambda node;
    \item new instances of subexpressions containing no free occurrences of the
    formal parameter will be constructed when they could safely and
    beneficially be shared (we discuss this point in more detail in Chapter
    15).
\end{numbered}

A more efficient alternative to this is \textit{compilation}, whereby we associate
with each lambda body a \textit{fixed sequence of instructions} which will construct an
instance of the lambda body. Then the operation of instantiating a lambda
body would consist simply of obeying the sequence of instructions associated
with the lambda body.

This instruction sequence can be constructed in advance by a compiler, and
contains implicitly the knowledge about the shape of the body and where the
formal parameter occurs. Hence we would expect the compiled code to run
much faster than the \ml{Instantiate} method, for just the same reasons that
compiled code runs faster than interpreted code in conventional languages. In
effect, all the tests in \ml{Instantiate} are made in advance by the compiler.
Furthermore, it turns out that compilation opens up many new avenues for
optimization, which offer considerable further efficiency increases.

Unfortunately, not all lambda abstractions are amenable to compilation in
this way. Consider, for example, the lambda abstraction

\begin{letalign}
    \tlb{x}(\tlb{y}$-$ y x)
\end{letalign}

When we apply the \ml{\tl x} abstraction to an argument, \ml{3} say, we instantiate its
body, thus creating a brand new lambda abstraction \ml{(\tlb{y}$-$ y 3)}.
Furthermore, each application of the \ml{\tl x} abstraction to a different argument
will create a new and different \ml{\tl y} abstraction, thus making a nonsense of our
hope to compile a single fixed code sequence for each lambda abstraction.

The problem is that \ml{x} occurs free in the body of the \ml{\tl y} abstraction, so that
we have to make a new instance of the \ml{\tl y} abstraction whenever \ml{x} is bound to a
new value by an application of the \ml{\tl x}  abstraction. In the case of lambda
abstractions which have no free variables there is no problem, and we can
compile a code sequence for it as outlined above.

One way around this problem would be to allow the code sequence to
access the values of the free variables in some way, thus parameterizing the
code sequence on the values of the free variables. This approach leads us to
the SECD machine [Landin, 1964; Henderson, 1980], in which the code
sequence for a lambda abstraction has access to an \textit{environment} which
contains values for each of the free variables, thus allowing a single code
sequence for each lambda abstraction. It is also the route followed by all
block-structured languages, in which the values of free variables are found by
looking in the appropriate stack frame.

In this book, however, we will study a totally different approach, called
\textit{supercombinator graph reduction}, which does not require the addition of an
environment to our model of graph reduction. The idea is to \textit{transform} the
program into an equivalent one in which all the lambda abstractions are
amenable to compilation. This transformation algorithm, which is called
\textit{lambda-lifting}, is of considerable interest in its own right, and we devote the
rest of this chapter and the next to it. After this, we spend a chapter discussing
an important optimization to lambda-lifting, called \textit{full laziness}, and Chapter
16 then digresses to describe an important alternative transformation, into \textit{SK
    combinators}. Finally, the bulk of the third part of the book (Chapters 18--21)
is spent in an extended discussion of how to compile the transformed program
into a linear instruction sequence, and the optimizations which this opens up.

\section{Solving the Problem of Free Variables}

In this section we outline our strategy for dealing with the problem of free
variables. We do so by using a modified form of \tb-reduction, in which we may
effectively perform several \tb-reductions at once.

Consider our current example
\begin{letalign}
    \tlb{x}\tlb{y}$-$ y x
\end{letalign}
Suppose we applied it to two arguments, thus:
\begin{letalign}
    (\tlb{x}\tlb{y}$-$ y x) 3 4
\end{letalign}
The lambda reducer described in Chapter~12 would proceed like this:
\begin{letalign}
    & (\tlb{x}\tlb{y}$-$ y x) 3 4 \\
    \reduction{} & (\tlb{y}$-$ y 3) 4 \\
    \reduction{} & $-$ 4 3
\end{letalign}

There is no reason, however, why we should not perform the \ml{\tl x} and \ml{\tl y}
reductions \textit{simultaneously}, thus:
\begin{letalign}
    & (\tlb{x}\tlb{y}$-$ y x) 3 4 \\
    \reduction{} & $-$ 4 3
\end{letalign}
This `multi-argument' reduction entails constructing an instance of the body
\ml{($-$ y x)} whilst substituting \ml{3} for free occurrences of \ml{x}, and \ml{4} for free
occurrences of \ml{y}. The following observations are crucial:

\begin{numbered}
    \item Much is gained by performing the reductions simultaneously. Firstly,
    doing so builds less intermediate structure in the heap, since the inter-
    mediate result of the \ml{\tl x} reduction is never constructed. Second (and
    more important), no problems are presented by the free occurrence of \ml{x}
    in the \ml{\tl y} abstraction.
    \item Nothing is lost by performing the \ml{\tl x} and \ml{\tl y} reductions simultaneously.
    The result of performing the \ml{\tl x} reduction alone is a \ml{\tl y} abstraction, and
    (assuming that we perform normal order reduction until WHNF is
    reached) no further work can be done on the \ml{\tl y} abstraction until it is given
    another argument.

    \phantom{xx}Hence we may as well wait until both arguments are present and then
    perform both reductions at once. This applies even if the application of the
    \ml{\tl x} abstraction to a single argument is shared.
\end{numbered}

\subsection{Supercombinators}

What sort of lambda abstractions are amenable to this sort of multi-argument
reduction? Simply lambda abstractions of the form \ml{(\tlb{x$_1$}\tlb{x$_2$}\ldots\tlb{x$_n$}E)}. This
motivates a new definition:

\definitionbox{\vspace{-1.5\baselineskip}}{%
    A \emph{supercombinator}, \ml{\$S}, of \emph{arity} $n$ is a lambda expression of the form
%    \vspace{-1.5\baselineskip}
    \begin{letalign}
        \tlb{x$_1$}\tlb{x$_2$}\ldots\tlb{x$_n$}E
    \end{letalign}
    where \ml{E} is not a lambda abstraction (this just ensures that all the `leading
    lambdas' are accounted for by \ml{x$_1$}\ldots\ml{x$_n$}) such that
    \vs
    \begin{compactenum}[(i)]
        \item \ml{\$S} has no free variables,
        \item any lambda abstraction in \ml{E} is a supercombinator,
        \item $n \geq 0$; that is, there need be no lambdas at all.
    \end{compactenum}
    \vs
    A \textit{supercombinator redex} consists of the application of a supercombinator
    to n arguments, where n is its arity. A \textit{supercombinator reduction} replaces
    a supercombinator redex by an instance of the supercombinator body with
    the arguments substituted for free occurrences of the corresponding
    formal parameter.
}

\noindent For example,
\begin{mlcoded}
    3 \\
    (+ 2 5) \\
    \tlb{x}x \\
    \tlb{x}+ x 1 \\
    \tlb{x}+ x x \\
    \tlb{x}\tlb{y}$-$ y x \\
    \tlb{f}f (\tlb{x}+ x x)
\end{mlcoded}
are all supercombinators, while the following are not:


\begin{tabular}{@{}l@{\quad}l@{}}
    \ml{\tlb{x}y }        & (\ml{y} occurs free) \\
    \ml{\tlb{y}$-$ y x}   & (\ml{x} occurs free) \\
    \ml{\tlb{f}f (\tlb{x}f x 2)} & (inner \ml{\tl x} abstraction is not a supercombinator, since \ml{f} \\
    & occurs free) \\
\end{tabular}
\vs

\subsubsection{Supercombinators of non-zero arity}

Supercombinators of non-zero arity (that is, having at least one \ml{\tl{}} at the front)
are important because they will be our \textit{unit of compilation}. Since they have no
free variables (clause (i)) we can compile a fixed code sequence for them.
Furthermore, clause (ii) ensures that any lambda abstractions in the body
have no free variables, and hence do not need to be copied when instantiating
the supercombinator body.

Such a supercombinator is somewhat analogous to a Pascal function which
takes several (value) parameters, which does not refer to any global variables,
and which has no side-effects.

\subsubsection{Supercombinators of arity zero and CAFs}

A supercombinator of arity zero (that is, having no \tl s at the front) is just a
constant expression (remember that it has no free variables). These super-
combinators are often called \textit{constant applicative forms} or CAFs. For
example,
\begin{mlcoded}
    3 \\
    + 4 5 \\
    + 3
\end{mlcoded}
are all CAFs. The last example makes the point that CAFs can still be
functions.

Since a CAF has no \ml{\tl}s at the front, they are never instantiated. Hence, no
code need be compiled for it, since a single instance of its graph can freely be
shared.

\subsubsection{Combinators}

A `supercombinator' sounds like a special sort of `combinator' and indeed this
is the case:

\definitionbox{
\vspace{-1.5\baselineskip}
}{
A \emph{combinator} is a lambda expression which contains no occurrences of a
free variable [Barendregt, 1984].
}

\noindent
A combinator is a `pure' function in the sense that the value of a combinator
applied to some arguments depends only on the values of the arguments, and
not on any free variables. The term `combinator' has a long pedigree [Curry
and Feys, 1958].

Thus some lambda expressions are combinators, and some combinators are
supercombinators.

\subsection{A Supercombinator-based Compilation Strategy}

If only all the lambda abstractions in our program were supercombinators!
Then it would be easy to compile them all, for the reasons mentioned in the
last section. Real programs, of course, have many lambda abstractions which
are not supercombinators, but it turns out to be relatively straightforward to
transform the program so that it contains only supercombinators. This will be
our strategy, and we embark on the transformation in Section~13.3.

For the sake of clarity we will often give names to supercombinators. These
names are entirely arbitrary, since the lambda abstractions are anonymous,
and we will normally begin them with a \ml{\$} to make them distinctive. Thus we
could write
\begin{letalign}
    \ml{\$XY} = \tlb{x}\tlb{y}$-$ y x
\end{letalign}
but to emphasize their special status further we will write the definition like
this:
\begin{letalign}
    \ml{\$XY} x y = $-$ y x
\end{letalign}
Our strategy is therefore to transform the lambda expression we wish to
compile into:
\begin{numbered}
    \item a set of supercombinator definitions, plus
    \item an expression to be evaluated.
\end{numbered}
To emphasize the inseparability of these two components we use a box, just as
we did in the case of Miranda programs (Section~3.3), thus:

\combinatorbox{
    Supercombinator definitions \\
    \phantom{XXXX}\ldots \\
    \phantom{XXXX}\ldots
}{
   Expression to be evaluated
}

\noindent
For example, we could represent the expression
\begin{letalign}
    (\tlb{x}\tlb{y}$-$ y x) 3 4
\end{letalign}
\noindent
as

\combinatorbox{
    \ml{\$XY x y = $-$ y x}
}{
    \ml{\$XY 3 4}
}

A crucial point in the definition of a supercombinator given above is that a
supercombinator reduction only takes place when all the arguments are
present. For example,
\begin{letalign}
    (\ml{\$XY} 3)
\end{letalign}
is not a supercombinator redex, and will not be reduced. We can therefore
regard the supercombinator definitions as a set of \textit{rewrite rules}. A reduction
consists of rewriting an expression which matches the left-hand side of a rule
with an instance of the corresponding right-hand side. Such systems are called
\textit{term rewrite systems} and have been much studied in their own right
[O'Donnell, 1977; Klop, 1980; Hoffman and O'Donnell, 1982].

\section[Transforming Lambda Abstractions]{Transforming Lambda Abstractions into Supercombinators}

To summarize our progress so far, we have seen that certain sorts of lambda
abstractions, the supercombinators, are particularly easy to compile. Our
implementation effort now breaks into two parts:

\begin{numbered}
    \item a translation algorithm which transforms all the lambda abstractions in
    the program into supercombinators;
    \item an implementation of supercombinator reduction.
\end{numbered}

First of all we consider how to transform lambda abstractions into super-
combinators. Here is an example program (in which neither lambda
abstraction is a supercombinator):

\combinatorbox{}{
    \ml{(\tlb{x}(\tlb{y}+ y x) x) 4}
}

Consider first the innermost lambda abstraction \ml{(\tlb{y}+ y x)}.

It has a free variable, \ml{x}, so it is not a supercombinator. However, a simple
transformation will make it into one:
\begin{quote}
make each free variable into an extra parameter (we sometimes call this
\textit{abstracting} the free variable).
\end{quote}
Thus we would transform
\begin{mlcoded}
    (\tlb{y}+ y x)
\end{mlcoded}
to
\begin{mlcoded}
    (\tlb{x}\tlb{y}+ y x) x
\end{mlcoded}
(This operation is simply \ml{\tb}-abstraction.) To see that these two expressions are
equivalent, just perform a \ml{\tb}-reduction on the second to get the first. To make
it slightly clearer we could perform an \ml{\ta}-conversion on the \ml{\tl x} abstraction to give
\begin{mlcoded}
    (\tlb{w}\tlb{y}+ y w) x
\end{mlcoded}
This clarifies the distinction between the two \ml{x}s which occurred in the previous
version. Now the lambda abstraction \ml{(\tlb{w}\tlb{y}+ y w)} is a supercombinator!
Performing this transformation on our original program gives

\combinatorbox{}{
    \ml{(\tlb{x}(\tlb{w}\tlb{y}+ y w) x x) 4}
}

Next we give the supercombinator a name, \ml{\$Y} say, like this

\combinatorbox{
    \ml{\$Y w y = + y w}
}{
    \ml{(\tlb{x}\$Y x x) 4}
}

Now we see that the \ml{\tl x} abstraction also fulfills the conditions for
supercombinatorhood, and we give it the name \ml{\$X}, thus

\combinatorbox{
    \ml{\$Y w y = + y w} \\
\indent    \ml{\$X x = \$Y x x}
}{
    \ml{\$X 4}
}

We can now execute our program by performing supercombinator
reductions:
\begin{letalign}
    & \ml{\$X 4} \\
    \reduction{} & \ml{\$Y 4 4} \\
    \reduction{} & \ml{+ 4 4} \\
    \reduction{} & \ml{8}
\end{letalign}
To review the algorithm so far:
\vs

\ml{UNTIL} there are no more lambda abstractions:
\begin{quote}
    \vspace{-0.5\baselineskip}
    \begin{compactenum}[(1)]
        \item Choose any lambda abstraction which has no inner lambda
        abstractions in its body.
        \item Take out all its free variables as extra parameters.
        \item Give an arbitrary name to the lambda abstraction (e.g.\ \ml{\$X34}).
        \item Replace the occurrence of the lambda abstraction by the name
        applied to the free variables.
        \item Compile the lambda abstraction and associate the name with the
        compiled code.
    \end{compactenum}
    \vspace{-0.5\baselineskip}
\end{quote}

\ml{END}
\vs

\noindent It is easy to see that we suffer an increase in the size of the program during this
transformation, but it is a price we pay willingly in exchange for the easier
reduction rules.

When we have completed the algorithm we arrive at a program of the form

\combinatorbox{
    {\ldots\ supercombinator definitions \ldots}
}{
    \ml{E}
}

But what about the expression \ml{E}? It must have no free variables, since it is
the top-level expression to be evaluated, so we can make it into a zero-
parameter supercombinator (a CAF) thus

\combinatorbox{
    {\ldots\ supercombinator definitions \ldots} \\
    \indent \ml{\$Prog = E}
}{
    \ml{\$Prog}
}

\noindent
thus completing the transformation of the program into supercombinators.

So far we have not made any explicit mention of recursion. This topic is so
important that we devote the whole of the next chapter to it.

Following Johnsson [1985], we call the transformation from lambda
expressions to supercombinators `lambda-lifting' since all the lambda
abstractions are lifted to the top level.

\subsection{Eliminating Redundant Parameters}

In this section and the next we will consider two simple optimizations to the
lambda-lifting algorithm. Consider the expression
\begin{mlcoded}
    \tlb{x}\tlb{y}$-$ y x
\end{mlcoded}
It is actually a supercombinator as it stands, but suppose we blindly applied
our algorithm as described above. First we choose the \ml{\tl y} abstraction, noting
that \ml{x} is free, and transform it to

\combinatorbox{
    \ml{\$Y x y = $-$ y x}
}{
    \ml{\tlb{x}\$Y x}
}

\noindent
(Here we have chosen to use \ml{x} instead of \ml{w} as the name of the extra parameter
to the \ml{\$Y} supercombinator. This choice is arbitrary, but we will normally
choose the same name as the free variable being abstracted.) Now dealing
with the \ml{\tl x} abstraction, we get

\combinatorbox{
    \ml{\$Y x y = $-$ y x} \\
    \ml{\$X x = \$Y x}
}{
    \ml{\$X}
}

It is clear that we can simplify the definition of \ml{\$X} to

\begin{letalign}
    \ml{\$X} = \ml{\$Y}
\end{letalign}

\noindent
(This is just \te-reduction, of course.) Having done this we see that \ml{\$X} itself is
redundant, and \ml{\$X} can be replaced wherever it occurs by \ml{\$Y}, giving

\combinatorbox{
    \ml{\$Y x y = $-$ y x}
}{
    \ml{\$Y}
}

So there are two optimizations to consider:

\begin{numbered}
    \item Remove redundant parameters from definitions by \te-reduction.
    \item Where this produces redundant definitions, eliminate them.
\end{numbered}

These optimizations together exploit supercombinators that appear naturally
in the original program, and sometimes catch other \te-reductions as well.

\textit{Caveat:} it turns out that, for more sophisticated implementations,
performing such \te-reductions is actually \textit{undesirable}, unless they succeed in
eliminating a definition, which is always desirable. For a full explanation of
this point, see Section~20.3.4.

\subsection{Parameter Ordering}

When we take out several free variables from a lambda abstraction as extra
parameters the order in which we put them seems rather arbitrary. For
example, consider the program

\combinatorbox{}{
    \ml{(\ldots} \\
    \phantom{XXXX} \ml{(\tlb{x}\tlb{z}+ y ($\ast$ x z))} \\
    \phantom{XXXX} \ml{\ldots)}
}

\noindent
where the `\ldots' stands for some expression enclosing the \ml{\tl x} abstraction. It
could be transformed to

\combinatorbox{
    \ml{\$S x y z = + y ($\ast$ x z)}
}{
    \ml{(\ldots} \\
    \phantom{XXXX} \ml{(\tlb{x}\$S x y)} \\
    \phantom{XXXX} \ml{\ldots)}
}

\noindent
or alternatively it could be transformed to

\combinatorbox{
    \ml{\$S y x z = + y ($\ast$ x z)}
}{
    \ml{(\ldots} \\
    \phantom{XXXX} \ml{(\tlb{x}\$S y x)} \\
    \phantom{XXXX} \ml{\ldots)}
}

\noindent
Both \ml{x} and \ml{y} are free, and it does not seem to matter which order we take them
out in. However, let us take the second possibility one stage further, by lifting
the \ml{\tl x} abstraction:

\combinatorbox{
    \ml{\$S y x z = + y ($\ast$ x z)} \\
    \indent \ml{\$T y x \phantom{z} = \$S y x}
}{
    \ml{(\ldots} \\
    \phantom{XXXX} \ml{(\$T y)} \\
    \phantom{XXXX} \ml{\ldots)}
}

\noindent
Now we can remove the redundant parameters from the definition of \ml{\$T}, and
eliminate the definition of \ml{\$T} altogether (since it is the same as \ml{\$S}). This
would not have been possible had we put the extra parameters \ml{x} and \ml{y} for \ml{\$S} in
the other order. Hence we should \textit{order the free variables}, with those bound at
inner levels coming last in the parameter list of the supercombinator.

This suggests that we could associate a \textit{lexical level-number} with each
lambda abstraction, so that the lexical level-number of a lambda abstraction is
defined to be one more than the number of textually enclosing lambdas (the
experienced reader will recognize these level-numbers as de Bruijn numbers
[de Bruijn, 1972]). For example, consider
\begin{mlcoded}
    (\tlb{x}\tlb{y}+ x ($\ast$ y y))
\end{mlcoded}
The \ml{\tl x} abstraction is at level 1, while the \ml{\tl y} abstraction is at level 2 (since it is
enclosed by a \ml{\tl x} abstraction).

The lexical level of a variable is now defined to be the lexical level of the
lambda abstraction which binds it. If the level of \ml{x} is less than \ml{y} we say that \ml{x} is
\textit{freer} than \ml{y}, since it is bound further out.

Constants (including built-in functions such as \ml{+}, and previously generated
supercombinators) can be regarded as being bound at the top level, and so
should be at level 0. There is, of course, no need to abstract out constants as
extra parameters during lambda-lifting.

To summarize:

\begin{numbered}
    \item The level-number of a \textit{lambda abstraction} is one more than the number
    of lambda abstractions which textually enclose it. If there is none, then
    its level-number is 1.
    \item The level-number of a \textit{variable} is the level-number of the lambda
    abstraction which binds it.
    \item The level-number of a \textit{constant} is 0.
\end{numbered}

It is simple to determine the lexical levels of all variables in a single
tree-walk over the expression. On the way down the tree the level-numbers of
the lambdas are recorded in a sort of environment, while on the way up the
level of each variable is computed, using the environment.

Now to maximize the chances of being able to apply \te-reduction we can
simply sort the extra parameters in increasing order of lexical level.

\section{Implementing a Supercombinator Program}

All the preceding work has shown how to compile our program into a set of
supercombinator definitions. What happens now? We have spoken of the
supercombinators being compiled in some way, but in fact there is a spectrum
of possible implementations:

\begin{numbered}
    \item We could keep the body of the supercombinator as a tree, and instantiate
    it using a function similar to \ml{Instantiate}. This is the supercombinator
    equivalent of the lambda reducer in the last chapter, and all the
    mechanisms described in the previous chapters concerning how to find
    the next redex, how to perform a reduction, indirection nodes, the stack
    and so on are still valid. The only change required is in the implemen-
    tation of \ml{Instantiate}. It is simplified because all lambda abstractions are
    known to be supercombinators (which have no free variables, and hence
    need never be copied), but is made more complicated because it has to
    substitute for several variables at once.

    \phantom{xx} We call this the template-instantiation implementation.

    \item We could keep the body of the supercombinator as a tree, but held in a
    contiguous block of store. Now the instantiation can be done with a
    modified block move, which can be implemented much more efficiently
    than a tree-walking instantiation. This idea is used by Keller [1985]. It is
    possible because supercombinators are constructed once and for all at
    compile-time, rather than being generated on the fly at run-time.

    \item We could compile the body to a linear sequence of instructions which will
    create an instance when executed. This is the idea behind the G-machine
    [Johnsson, 1984], which we discuss in Chapters 18--21. This is faster still,
    and also opens the way to many further optimizations, as we shall see.
    \end{numbered}

    \noindent
    The fundamental point is that all we can do with a supercombinator is to apply
    it, and hence we are free to choose a representation for the supercombinator
    that makes this operation efficient.

    \section*{References}

    \begin{references}

    \item Barendregt, H.P. 1984. \textit{The Lambda Calculus: its syntax and semantics}, 2nd edition, p.
    24. North-Holland.

    \item Curry, H.B., and Feys, R. 1958. \textit{Combinatory Logic}, Vol. I. North-Holland.

    \item De Bruijn, N.G. 1972. Lambda calculus notation with nameless dummies.
    \textit{Indagationes Mathematicae}. Vol. 34, pp. 381--92.

    \item Henderson, P. 1980. \textit{Functional Programming: application and implementation}.
    Prentice-Hall.

    \item Hoffman, C.M., and O'Donnell, M.J. 1982. Programming with equations. \textit{ACM
        TOPLAS}. Vol. 4, no. 1, pp. 83--112.

    \item Hughes, R.J.M. 1984. The design and implementation of programming languages.
    PhD thesis, PRG-40, Programming Research Group, Oxford. September.

    \item Johnsson, T. 1984. Efficient compilation of lazy evaluation. In \textit{Proceeding of the ACM
        Symposium on Compiler Construction, Montreal}, pp. 58--69. June.

    \item Johnsson, T. 1985. Lambda lifting: transforming programs to recursive equations. In
    \textit{Conference on Functional Programming Languages and Computer Architecture,
        Nancy}. Jouannaud (editor). LNCS 201. Springer Verlag.

    \item Keller, R.M. 1985. Distributed graph reduction from first principles. Department of
    Computer Science, University of Utah.

    \item Klop, J.W. 1980. Combinatory reduction systems. PhD thesis, Mathematisch
    Centrum, Amsterdam.

    \item Landin, P.J. 1964. The mechanical evaluation of expressions. \textit{Computer Journal}. Vol.
    6, pp. 308--20.

    \item O'Donnell, M.J. 1977. \textit{Computing in Systems Described by Equations}. LNCS 58.
    Springer Verlag.

    \end{references}